\documentclass[11pt]{article}
\usepackage[utf8]{inputenc}
\usepackage[spanish]{babel}
\usepackage[paperwidth=4.7in, paperheight=9in, margin=0.35in]{geometry}
\usepackage{hyperref}
\usepackage{booktabs}
\usepackage{tabularx}
\usepackage{xcolor}
\usepackage{enumitem}
\usepackage{titlesec}
\usepackage{microtype}
\usepackage{tikz}
\usetikzlibrary{positioning, calc, arrows.meta, shapes.geometric}

% Colores premium estilo ELITE FTNS
\definecolor{neonYellow}{HTML}{F3CA4C}
\definecolor{darkBg}{HTML}{121212}
\definecolor{lightGray}{HTML}{F9F9F9}
\definecolor{linkBlue}{HTML}{1A73E8}

\hypersetup{
    colorlinks=true,
    linkcolor=linkBlue,
    urlcolor=linkBlue,
    pdftitle={Documento de Arquitectura ELITE FTNS}
}

% Ajuste de fuentes y espaciados para lectura cómoda en celular
\renewcommand{\familydefault}{\sfdefault}
\titleformat{\section}{\large\bfseries\color{darkBg}}{\thesection}{1em}{}
\titleformat{\subsection}{\normalsize\bfseries\color{darkBg}}{\thesubsection}{1em}{}
\titlespacing*{\section}{0pt}{1.5ex plus 1ex minus .2ex}{1ex plus .2ex}
\titlespacing*{\subsection}{0pt}{1.2ex plus 1ex minus .2ex}{0.8ex plus .2ex}

\setlist[itemize]{leftmargin=12pt, noitemsep, topsep=2pt}
\setlist[enumerate]{leftmargin=12pt, noitemsep, topsep=2pt}

\title{\textbf{\large ELITE FTNS \\ Documento de Arquitectura y Guía de Integración}}
\author{\small Antigravity Co-Pilot}
\date{\footnotesize \today}

\begin{document}

\maketitle

\section{Introducción}
Este documento consolida la arquitectura técnica, las especificaciones de frontend, backend y el flujo de integración de funcionalidades core de la plataforma \textbf{ELITE FTNS (v2.0.3)}. Está diseñado para su consulta en dispositivos móviles durante viajes de investigación.

---

\section{Arquitectura y Backend}
El backend está implementado en **FastAPI** y estructurado bajo un modelo de datos multi-tenant aislado a nivel de archivos físicos de base de datos.

\subsection{Estructura Multi-Tenant}
El aislamiento de entrenadores se maneja mediante:
\begin{enumerate}
    \item \textbf{Base de Datos Maestra (\texttt{master.db}):} Administra las credenciales y el estado de suscripción de los entrenadores (tabla \texttt{trainers}).
    \item \textbf{Bases de Datos de Tenants (\texttt{trainer\_*.db}):} Ubicadas en el directorio \texttt{database/tenants/}. Cada entrenador cuenta con su propio archivo SQLite físico, aislando completamente la información de sus clientes, rutinas, valoraciones y dietas.
\end{enumerate}

\begin{figure}[h]
\centering
\begin{tikzpicture}[
    node distance=0.6cm,
    box/.style={draw=darkBg, thick, rectangle, rounded corners, fill=lightGray, text width=3.3in, align=center, minimum height=0.6cm, font=\footnotesize},
    arrow/.style={-{Latex}, thick, darkBg}
]
    \node (master) [box] {\textbf{Master DB} (\texttt{master.db}) \\ Tabla: \texttt{trainers}};
    \node (tenantDir) [box, below=of master] {\textbf{Directorio de Tenants} (\texttt{database/tenants/})};
    \node (tenant1) [box, below=0.4cm of tenantDir, text width=1.5in, xshift=-0.85in] {\textbf{Tenant Admin} \\ \texttt{trainer\_admin.db}};
    \node (tenant2) [box, below=0.4cm of tenantDir, text width=1.5in, xshift=0.85in] {\textbf{Tenant Carlos} \\ \texttt{trainer\_carlos.db}};
    
    \draw [arrow] (master) -- (tenantDir);
    \draw [arrow] (tenantDir.south) -- ++(0,-0.2) -| (tenant1.north);
    \draw [arrow] (tenantDir.south) -- ++(0,-0.2) -| (tenant2.north);
\end{tikzpicture}
\caption{Aislamiento Físico de Datos (Multi-Tenant).}
\end{figure}

\subsection{Modo WAL de SQLite}
Para permitir lecturas y escrituras concurrentes (crítico cuando múltiples clientes registran datos de entrenamiento en vivo), las bases de datos de inquilinos se inicializan con la instrucción:
\begin{quote}\footnotesize
    \texttt{PRAGMA journal\_mode=WAL;}
\end{quote}
Esto evita bloqueos de base de datos durante escrituras de bitácora y mensajería WebSocket.

\subsection{Lógica de Autosemillado para Nuevos Tenants}
Cuando se registra un nuevo entrenador, el backend copia automáticamente la biblioteca global de alimentos (\texttt{food\_library}) y el catálogo de ejercicios base (\texttt{exercises}) desde la base de datos maestra a su nuevo archivo tenant, asegurando un entorno de trabajo funcional desde el primer segundo.

\subsection{Lógica de Persistencia de la Bitácora}
En el endpoint de guardado diario de datos del cliente, el servidor utiliza una consulta con \texttt{sqlite3.Row} para recuperar el registro diario previo (si existe), realizar un \textbf{merge} de los datos parciales (hidratación, peso, checklists) y guardar la estructura combinada. Esto previene que una actualización parcial de hidratación sobreescriba o borre el checklist de ejercicios completados del mismo día.

\subsection{API de Endpoints Principales}
\begin{itemize}
    \item \texttt{POST /api/auth/login:} Autenticación segura con JWT.
    \item \texttt{GET/POST /api/clients:} Gestión de clientes por tenant.
    \item \texttt{GET/POST/PUT /api/assessments:} Valoraciones físicas (cálculo de Faulkner, FFMI, etc.).
    \item \texttt{GET/POST /api/assessment\_config:} Configuración dinámica de campos.
    \item \texttt{PUT /api/clients/<id>/routine:} Cambiar rutina asignada.
    \item \texttt{WS /ws/chat:} Canal dúplex en tiempo real para chat.
\end{itemize}

---

\section{Diseño y Frontend}
El frontend está desarrollado con tecnologías web estándar (HTML, CSS vanilla y JavaScript modular), empaquetado como una **PWA (Progressive Web App)** con soporte sin conexión para el cliente.

\subsection{Identidad Visual: Energía Neón}
Se define por un diseño estilo \textbf{dark glassmorphism}:
\begin{itemize}
    \item \textbf{Fondo oscuro} con acentos de color amarillo neón (\texttt{\#F3CA4C}).
    \item \textbf{Efecto Glassmorphism:} Tarjetas y contenedores con un fondo negro translúcido y desenfoque por software:
    \begin{quote}\footnotesize
        \texttt{background: rgba(18, 18, 18, 0.7); \\ backdrop-filter: blur(10px);}
    \end{quote}
\end{itemize}

\subsection{PWA (Bitácora Móvil)}
Mediante un \texttt{manifest.json} y un \texttt{service-worker.js}, el portal de cliente se puede agregar a la pantalla de inicio de celulares Android e iOS, ofreciendo navegación de pantalla completa y caché local.

\subsection{Responsividad y Breakpoints}
*   El breakpoint global se fijó en \textbf{\texttt{1024px}} (en lugar de \texttt{768px}). Esto asegura que las vistas en modo horizontal (Landscape) de tabletas y móviles mantengan el diseño de aplicación nativa (con la barra inferior de navegación en lugar de la barra lateral de escritorio).
*   Las tablas de datos fueron optimizadas reduciendo el padding a \texttt{4px}, disminuyendo el tamaño de la tipografía a \texttt{11px} e incorporando un scroll horizontal automático (\texttt{overflow-x: auto}) para evitar deformaciones en pantallas pequeñas.

\subsection{Resolución de Capas (z-index)}
El panel deslizable de la ficha del cliente en móvil tiene un \texttt{z-index: 10000}. Para evitar que este panel tape a las ventanas emergentes (crear cliente, cambiar rutina, etc.), se estableció globalmente en \texttt{style.css} una regla de z-index superior para todos los modales:
\begin{quote}\footnotesize
    \texttt{.modal-container \{ z-index: 20000 !important; \}}
\end{quote}

---

\section{Integración de Funcionalidades Core}

\subsection{Checklists en Tiempo Real}
El Portal del Cliente incluye checkboxes al lado de cada ejercicio y alimento.
1.  **Sincronización:** Al pulsar un checkbox, JavaScript envía el estado del ejercicio/comida de inmediato a través de una petición de guardado.
2.  **Visualización del Entrenador:** Estos progresos se reflejan de inmediato en el Portal del Entrenador al abrir el día correspondiente en el calendario mensual de trazabilidad.

\subsection{Flujo de ``Cambiar Rutina''}
Se eliminaron las opciones de asignación y desasignación separadas para evitar redundancia:
1.  El entrenador abre el modal unificado de **``Cambiar Rutina''**.
2.  Un buscador dinámico en JavaScript filtra en vivo las tarjetas interactivas de las rutinas globales por título.
3.  Al presionar ``Guardar Cambios'', el backend desvincula la rutina anterior y asocia la nueva de manera atómica.

\subsection{Chat en Tiempo Real por WebSockets}
El chat funciona con un canal WebSocket que orquesta la entrega instantánea:

\begin{figure}[h]
\centering
\begin{tikzpicture}[
    node distance=1.2cm,
    every node/.style={font=\scriptsize},
    line/.style={draw, thick, dashed, darkBg},
    actor/.style={draw=darkBg, thick, rectangle, rounded corners, fill=lightGray, text width=0.9in, align=center, minimum height=0.5cm},
    msg/.style={-{Latex}, thick, darkBg}
]
    \node (c) [actor] {\textbf{Cliente}};
    \node (s) [actor, right=1.0in of c] {\textbf{FastAPI WS}};
    \node (db) [actor, right=1.0in of s] {\textbf{SQLite Tenant}};
    
    \draw [line] (c.south) -- ++(0,-4.5);
    \draw [line] (s.south) -- ++(0,-4.5);
    \draw [line] (db.south) -- ++(0,-4.5);
    
    \draw [msg] ($(c.south)+(0,-0.6)$) -- node[above] {1. Mensaje WS} ($(s.south)+(0,-0.6)$);
    \draw [msg] ($(s.south)+(0,-1.6)$) -- node[above] {2. INSERT SQL} ($(db.south)+(0,-1.6)$);
    \draw [msg, dashed] ($(db.south)+(0,-2.6)$) -- node[above] {3. ID generado} ($(s.south)+(0,-2.6)$);
    \draw [msg, dashed] ($(s.south)+(0,-3.6)$) -- node[above] {4. Sent Receipt} ($(c.south)+(0,-3.6)$);
\end{tikzpicture}
\caption{Diagrama de Flujo del Chat de Mensajes.}
\end{figure}

El sistema implementa:
*   **Burbujas translúcidas:** Fondo oscuro con desenfoque de fondo y acentos azul/cian en mensajes enviados.
*   **Ticks de Confirmación:** Un tick al guardarse en el servidor, doble tick al entregarse (receptor online), y doble tick coloreado al leerse.

\subsection{Calendario de Trazabilidad Diaria}
Se solucionó un bug de colisión de IDs en el DOM eliminando el bloque inactivo de inyección y centralizando la carga únicamente en el modal emergente (\texttt{dayDetailModal}). Al pulsar un día en el calendario, se abre el modal mostrando peso, hidratación, notas del día y el checklist completado.

\end{document}
